\section{Discussion and Limitations}\label{sec:limitation}

\tool provides an uncertainty-aware framework for probabilistic fault diagnosis, but several limitations remain. First, the current formulation focuses on processor faults and treats edge reliability as a confidence weight for syndrome evidence. A more general model could explicitly diagnose both faulty nodes and faulty links as two coupled latent variables. Second, the training data are generated from simulated probabilistic PMC syndromes. Although simulation is standard in system-level diagnosis, future work should validate the method on traces from real high-performance computing or data-center systems. Third, the current posterior approximation is learned by a neural model and does not provide the same type of worst-case diagnosability guarantee as classical graph-theoretic diagnosis. Combining calibrated learning-based diagnosis with formal diagnosability bounds is an important future direction.

Another limitation is that calibration can degrade under severe out-of-distribution shifts. Parameter-randomized training improves robustness, but it cannot cover all possible operating conditions. In practice, \tool should be used together with monitoring mechanisms that detect when the test environment has shifted beyond the training range. Finally, top-$k$ suspicious-node ranking is useful for maintenance, but the optimal value of $k$ may depend on operational cost, repair budget, and service-level requirements.